\documentclass[11pt]{article}

\usepackage[margin=1in]{geometry}
\usepackage{amsmath,amssymb}
\usepackage{array}
\usepackage{booktabs}
\usepackage{longtable}
\usepackage{xcolor}
\usepackage{enumitem}
\usepackage{hyperref}

\hypersetup{
  colorlinks=true,
  linkcolor=blue!60!black,
  urlcolor=blue!60!black,
  citecolor=blue!60!black
}

\setlist[itemize]{leftmargin=1.5em}
\setlist[enumerate]{leftmargin=1.7em}

\newcommand{\cert}[1]{\path{#1}}
\newcommand{\level}[1]{\textbf{\textsf{#1}}}
\newcommand{\GL}{\operatorname{GL}}
\newcommand{\Sym}{\operatorname{Sym}}
\newcommand{\Hilb}{\operatorname{Hilb}}
\newcommand{\wt}{\operatorname{wt}}

\title{Technical report on the second literal assertion in Lee's Lemma 19}
\author{Haiman counterexample audit project}
\date{2026-06-27}

\begin{document}
\maketitle

\begin{abstract}
This report expands the short section
\href{https://nasqret.github.io/haiman-counterexample-audit/chapters/lemma19.html#a-second-literal-assertion-fails}{``A second literal assertion fails''}
in the project documentation.  The purpose is narrow: explain exactly what was
computed, which representation-theoretic facts are being used, why no Borel
choice is needed for this particular negative test, and what remains open.

The current audit proves that one printed sentence in Lee's Lemma 19 is false
if read literally: not every \(89\times 89\) minor belongs to one of the 15
degree-89 predecessor modules.  This is a serious gap in the written proof, but
it is not by itself a proof that the claimed degree-90 generator is nonminimal,
nor a proof that Lee's theorem is false.
\end{abstract}

\section{Executive summary}

The source paper claims that a degree-90 Schur module
\[
  S_\nu W,\qquad
  \nu=(133,130,126,122,119,60,60,60),
\]
occurs among the minimal generators of \(J_8\), the ideal of the principal
component \(P_8 \subset \Hilb^9(\mathbb C^8)\) near \(\mathfrak m^2\).  Lemma 19
is the computational step intended to prove this.

The audit has reconstructed Lee's \(90\times115\) matrix \(M\) and certified a
nonzero \(90\times90\) minor.  The section studied here concerns a different
sentence on the degree-89 side.  The paper says, in effect, that each
\(89\times89\) minor belongs to one of 15 modules obtained by a
Littlewood--Richardson search.  The audit produced an explicit counterexample
to that literal sentence.

\begin{center}
\begin{tabular}{p{0.22\linewidth}p{0.68\linewidth}}
\toprule
Evidence level & Statement \\
\midrule
\level{Computation} &
Deleting row 0 and the first pivot column from the certified full-rank
\(90\times90\) pivot minor gives a nonzero \(89\times89\) determinant:
\[
421057 \pmod {1000003}.
\]
\\
\level{Computation} &
The same minor is torus-weight homogeneous of weight
\[
(60,59,59,140,123,126,119,115),
\]
whose dominant reordering is
\[
(140,126,123,119,115,60,59,59).
\]
\\
\level{Computation} &
The 15 certified degree-89 predecessor partitions all have first part at most
132. \\
\level{Theorem} &
Every weight of the irreducible \(\GL(W)\)-module \(S_\lambda W\) has dominant
representative dominated by \(\lambda\).  In particular, its first part is at
most \(\lambda_1\). \\
\level{Conclusion} &
The explicit nonzero \(89\times89\) minor is not a weight vector in any of
the 15 listed degree-89 modules.  Hence the literal sentence ``each
\(89\times89\) minor belongs to one of the 15 modules'' is false as written. \\
\bottomrule
\end{tabular}
\end{center}

\section{Why torus weights are enough here}

The user's concern is exactly the right one: the published webpage section is
short and mentions torus weights, while highest-weight statements usually
involve a Borel subgroup.  The distinction is as follows.

\subsection{Positive membership needs more than a weight}

To prove that a vector generates or lies in a particular copy of
\(S_\lambda W\), a torus weight alone is not sufficient.  One normally also
checks a highest-weight condition for a chosen Borel, or computes the
projection to the desired isotypic component.  This is why the audit does not
claim that the weight-compatible \(90\times90\) minor proves Lee's
Schur-module assertion.

Indeed, the audit found a nonzero \(90\times90\) minor of the claimed torus
weight
\[
(60,60,60,133,130,126,122,119),
\]
with determinant \(300834 \pmod {1000003}\).  But for the Borel order
\[
4<5<6<7<8<1<2<3,
\]
five simple raising derivatives are nonzero:
\[
\begin{array}{c|ccccc}
 \text{operator} & E_{45} & E_{56} & E_{67} & E_{78} & E_{23} \\
\hline
 \text{residue mod }1000003 & 685026 & 176188 & 140239 & 78485 & 605583 .
\end{array}
\]
Thus that selected determinant is not itself a highest-weight vector.  This is
a negative computation about one determinant only; it does not rule out a
different linear combination or a nonzero projection.

\subsection{Negative membership can be disproved by a missing weight}

The \(89\times89\) statement is different.  We are not trying to prove positive
membership in one of the 15 modules.  We are disproving membership by showing
that the vector has a torus weight that cannot occur in any of them.

Let \(T\subset\GL(W)\) be the diagonal torus.  Any finite-dimensional rational
\(\GL(W)\)-module decomposes into \(T\)-weight spaces.  An irreducible
\(S_\lambda W\) has only weights whose dominant representative is dominated by
\(\lambda\).  This is a standard highest-weight theorem, but the test itself
does not require choosing a Borel after the list of highest weights is fixed.

Therefore, if a homogeneous determinant \(f\) has torus weight \(\alpha\), and
the dominant reordering \(\alpha^+\) is not dominated by any of the 15
candidate highest weights \(\lambda\), then \(f\) cannot be an element of their
direct sum.  It cannot be hidden there under another Borel convention; the
torus weight is already impossible.

\section{The reconstructed matrix and the certified minor}

Lee's construction isolates the 45 variables
\[
 A=\{p_{r,st}:1\le r\le3,\;4\le s\le t\le8\}.
\]
After centering the complementary variables, the selected equations have the
form
\[
 C=M(A)z,
\]
where \(M\) is a \(90\times115\) matrix with entries \(0,\pm p_{r,st}\).

The matrix reconstruction is not a transcription of entries.  It is generated
from the printed index formulas.  The canonical sparse manifest has:
\[
\begin{array}{rl}
\text{shape} &= 90\times115,\\
\text{coefficient variables} &= 45,\\
\text{nonzero entries} &= 1410.
\end{array}
\]
The payload SHA-256 is
\[
\text{\scriptsize\ttfamily dc4bb49eaab862c6f11ce83b9129fd9c16adc6b6d4e591fc4be8ea972d9176ac}.
\]

At the recorded specialization modulo \(p=1000003\), the matrix has rank 90.
The pivot minor has determinant
\[
  970351 \pmod {1000003}.
\]
This proves that the corresponding integer determinant polynomial is not zero
over characteristic zero.

For the degree-89 test, the audit deletes:
\[
 \text{row }0,\qquad \text{pivot-column position }0
 \quad(\text{global column }0)
\]
from the same certified pivot minor.  The resulting \(89\times89\) determinant
is
\[
  421057 \pmod {1000003}.
\]
Thus this \(89\times89\) determinant is also a nonzero polynomial in
characteristic zero.

\section{Weight calculation}

The coordinate \(p_{r,st}\) has torus weight
\[
 \wt(p_{r,st})=(1,\ldots,1)-e_r+e_s+e_t.
\]
The matrix entries are compatible with row and column weights: an entry in row
\(\rho\), column \(\gamma\), and coefficient variable \(a\) satisfies
\[
 \wt(a)=\wt(\rho)-\wt(\gamma).
\]
Consequently every nonzero monomial in a fixed minor determinant has the same
total torus weight:
\[
 \wt(\det M_{R,C})
 =
 \sum_{\rho\in R}\wt(\rho)-\sum_{\gamma\in C}\wt(\gamma).
\]

For the certified \(89\times89\) minor this gives
\[
\alpha=(60,59,59,140,123,126,119,115).
\]
Reordering into the dominant chamber gives
\[
\alpha^+=(140,126,123,119,115,60,59,59).
\]

\section{The 15 predecessor partitions}

The audit reconstructs the Littlewood--Richardson search appearing in Lee's
Lemma 19.  Let
\[
\mu=(3,1,1,1,1,1,1,0),\qquad
\nu=(133,130,126,122,119,60,60,60).
\]
The certificate records:
\begin{itemize}
  \item 102 partitions \(\lambda\) with \(c^\nu_{\lambda,\mu}>0\);
  \item 87 excluded by tensor-product dominance for the 89th tensor power;
  \item 15 remaining candidates, each with an explicit 88-step nonzero
  Littlewood--Richardson chain proving tensor occurrence.
\end{itemize}

The 15 endpoint partitions are:

\begin{center}
\begin{longtable}{r>{\(\displaystyle}l<{\)}}
\toprule
\# & \lambda\\
\midrule
1 & (130,129,125,121,118,60,59,59)\\
2 & (131,128,125,121,118,60,59,59)\\
3 & (131,129,124,121,118,60,59,59)\\
4 & (131,129,125,120,118,60,59,59)\\
5 & (131,129,125,121,117,60,59,59)\\
6 & (132,127,125,121,118,60,59,59)\\
7 & (132,128,124,121,118,60,59,59)\\
8 & (132,128,125,120,118,60,59,59)\\
9 & (132,128,125,121,117,60,59,59)\\
10 & (132,129,123,121,118,60,59,59)\\
11 & (132,129,124,120,118,60,59,59)\\
12 & (132,129,124,121,117,60,59,59)\\
13 & (132,129,125,119,118,60,59,59)\\
14 & (132,129,125,120,117,60,59,59)\\
15 & (132,129,125,121,116,60,59,59)\\
\bottomrule
\end{longtable}
\end{center}

Every one has \(\lambda_1\le132\).

\section{The actual contradiction to the literal sentence}

For a weight \(\alpha\) to occur in \(S_\lambda W\), its dominant reordering
\(\alpha^+\) must be dominated by \(\lambda\).  Dominance begins with the
first prefix inequality:
\[
 \alpha^+_1 \le \lambda_1.
\]
For the certified \(89\times89\) minor:
\[
 \alpha^+_1=140.
\]
For every one of the 15 predecessor candidates:
\[
 \lambda_1\le132.
\]
Therefore
\[
 \alpha^+\nleq \lambda
\]
for all 15 candidates.  The determinant cannot be a weight vector in any of
the 15 modules.  Since the determinant is nonzero, this is a concrete
counterexample to the literal reading of the sentence that each
\(89\times89\) minor belongs to one of those modules.

\section{What this does and does not prove}

\subsection{What is proved}

\begin{enumerate}
  \item The \(89\times89\) determinant specified in
  \cert{results/certificates/lemma19_89_minor_weight_audit.json} is
  nonzero in characteristic zero.
  \item Its torus weight is
  \((60,59,59,140,123,126,119,115)\).
  \item This weight cannot occur in any of the 15 certified predecessor
  irreducibles.
  \item Therefore the printed universal statement about all \(89\times89\)
  minors is false if read literally.
\end{enumerate}

\subsection{What is not proved}

\begin{enumerate}
  \item This does not prove that Lee's claimed degree-90 generator is
  nonminimal.
  \item This does not prove that \(S_\nu W\) is absent from the span of all
  \(90\times90\) maximal minors.
  \item This does not prove that the theorem of the paper is false.
  \item This does not rule out a corrected interpretation such as: ``each
  relevant \(89\times89\) component that can contribute to \(S_\nu W\) lies in
  one of these 15 modules.''
\end{enumerate}

The last point is important.  The displayed \(89\times89\) minor has dominant
first part 140, so it is not even contained in \(\nu\), whose first part is
133.  Hence this particular \(89\times89\) module is not an obvious source for
the target \(S_\nu W\) after multiplication by the degree-one coordinate
representation.  It refutes the paper's broad wording, but it does not by
itself exhibit a lower-degree generator that produces the target degree-90
copy.

\section{How it could still affect minimality}

Lee's intended minimality argument must show that the target degree-90 copy of
\(S_\nu W\) is not generated by degree-89 elements of \(J_8\).  A complete
certificate should identify the multiplication map
\[
 R_1\otimes (J_8)_{89}\longrightarrow (J_8)_{90}
\]
and prove that the relevant \(S_\nu W\)-projection is not in the image.

The current \(89\times89\) counterexample matters because it shows that the
published bookkeeping of degree-89 minors is not literally complete.  This
creates two possible outcomes:

\begin{itemize}
  \item \textbf{Repair path.}  The author only needed the 15 modules that can
  map to \(S_\nu W\), and the extra \(89\times89\) minors are irrelevant to the
  target projection.  Then the paper's wording is wrong, but the construction
  may still be repairable.
  \item \textbf{Failure path.}  The missing degree-89 bookkeeping hides a
  nonzero image into the claimed degree-90 constituent, making the purported
  degree-90 generator nonminimal.  Proving this requires an explicit
  projection or multiplication certificate, not just the weight obstruction in
  this report.
\end{itemize}

\section{Evidence table}

\begin{center}
\small
\begin{tabular}{>{\raggedright\arraybackslash}p{0.21\linewidth}
                >{\raggedright\arraybackslash}p{0.34\linewidth}
                >{\raggedright\arraybackslash}p{0.33\linewidth}}
\toprule
Item & Source of evidence & Status \\
\midrule
Matrix shape \(90\times115\) and 1410 sparse entries &
SageMath generator plus Rust, Singular, Oscar, Magma verifiers &
\level{Computation/certificate} \\
\midrule
Full-rank \(90\times90\) specialization &
\cert{lemma19_nonzero_minor.json}; determinant \(970351\) modulo \(1000003\) &
\level{Computation/certificate} \\
\midrule
Nonzero \(89\times89\) minor &
\cert{lemma19_89_minor_weight_audit.json}; determinant \(421057\) modulo
\(1000003\) &
\level{Computation/certificate} \\
\midrule
Weight formula for coordinates and determinants &
Torus action on \(p_{r,st}\); additive row/column factorization &
\level{Theorem plus checked bookkeeping} \\
\midrule
List of 15 predecessor partitions &
Sage LR enumeration plus independent Rust LR-tableau verifier &
\level{Computation/certificate} \\
\midrule
Weight exclusion from the 15 modules &
Highest-weight theory: weights of \(S_\lambda W\) are dominated by \(\lambda\) &
\level{Theorem} \\
\midrule
Degree-90 minimality &
Requires the \(S_\nu W\)-projection of \(R_1(J_8)_{89}\to(J_8)_{90}\) &
\level{Open} \\
\midrule
Lee's theorem and Haiman counterexample &
Requires Lemma 19 plus downstream regularity and geometry audit &
\level{Open} \\
\bottomrule
\end{tabular}
\end{center}

\section{Reproduction commands}

The relevant local commands are:
\begin{verbatim}
./scripts/validate-all.sh

python3 artifacts/common/analyze_minor_weight.py

cargo run --quiet --locked \
  --manifest-path artifacts/rust/Cargo.toml
\end{verbatim}

The important certificate files are:
\begin{itemize}
  \item \cert{results/certificates/lemma19_89_minor_weight_audit.json};
  \item \cert{results/certificates/lemma19_claimed_weight_raising_audit.json};
  \item \cert{results/certificates/lemma19_predecessor_partitions.json};
  \item \cert{results/certificates/lemma19_nonzero_minor.json}.
\end{itemize}

\section{Recommended next checks}

\begin{enumerate}
  \item Replace the phrase ``all \(89\times89\) minors'' by the precise
  representation-theoretic subspace that can contribute to \(S_\nu W\), if such
  a repair exists.
  \item Compute the \(S_\nu W\)-isotypic projection of the maximal-minor span.
  \item Compute the \(S_\nu W\)-projection of the multiplication image from
  degree 89.
  \item Decide whether the target degree-90 copy is outside that image
  (minimality) or inside that image (nonminimality).
  \item Only after that, return to the geometric and regularity steps of the
  paper.
\end{enumerate}

\section{Current verdict}

The current verdict remains \textbf{open audit}.  The second literal assertion
in Lemma 19 is refuted as written, and this is a natural place to investigate
the possible nonminimality of the claimed degree-90 generator.  However, the
available evidence does not yet prove nonminimality.  The next decisive object
is not another isolated torus weight, but the actual \(S_\nu W\)-projection of
the degree-90 minors and the multiplication image from degree 89.

\end{document}
